% !TeX root = ../papers/nc-p3.tex
\section{無印の欠色窓の二直線配置と三角PBW極限}
\label{sec:unmarked-klein-pbw}
\begin{theorem}\label{thm:unmarked-klein-pbw}
四位置のラベルを $K=(\mathbf Z/2)^2$ とし、総次数零、両隣接切断の色 $0$ は欠け、
三つの他色は全て階数2とする。四位置の全てで、$i\ne j$ の座標点対 $(i,j)$ は
辺の点対応に属さないとする。四文字が相異なる成分の消失は任意に許す。
\begin{enumerate}
\item $w_{ijij}=0$ が相異なるラベルの少なくとも一対で起これば、窓は常に標準14641型quadratic Koszul左右AS正則である。
ある完全マッチングの両辺の反復対が欠ける場合には、元の座標文字を並べ替えるだけで二次PBW表示となる。
それ以外も同じPBW窓への基底変更の一径数極限を持つ。
\item 全ての $w_{ijij}$ が非零なら、どの座標を星に取っても既出の四元群交換子表示が得られ、
定理\ref{thm:marked-klein-mixed-torus}から定理\ref{thm:marked-klein-zero-bracket-torus}までの全点判定が適用できる。
\end{enumerate}
従ってこの無印領域で全ての正則性が決まり、前節に残した完全マッチング、道、四サイクルの
全係数非零の場合はいずれも全点で正則である。
\end{theorem}
\begin{proof}[二次切断の行と列を同じ四つの直線で読む]
$M_b(i,k)=w_{i,i+b,k,k+b}$ と置く。四位置の点対についての仮定は、
$M_b$ と巡回して得られる他切断の各色行列に零行も零列もないことと同値である。
階数分解 $M_b(i,k)=l_i(v_k)$ を二次元空間で取り、非零な各 $l_i$ の核を
$p_i\in\mathbb P^1$ とする。欠色0から $M_b(i,i+b)=0$ であり、
$v_{i+b}$ は非零なので $kv_{i+b}=p_i$ となる。従って
\[
                 M_b(i,k)=0\quad\Longleftrightarrow\quad p_i=p_{k+b}.
\]
特に左辺の零・非零は $(i,k)\mapsto(k+b,i+b)$ で保たれる。二切断で使うと
四位置の反転 $(03)(12)$ と $(01)(23)$ がともに台を保つ。
四文字が相異なる成分 $w_{i,i+b,i+b+e,i+e}$ では、これら二つの反転が生成する四元群は
同じ $(b,e)$ の四つの $i$ に推移的に作用する。従ってこの四項の非零・零は一斉である。
反復対でも $w_{ijij}$ と $w_{jiji}$ の消失は同値である。

$b\ne e$ で四項が零なら $p_i=p_{i+e}$ が全 $i$ に成立する。
残る二組の直線は相異なる。そうでなければ $M_b$ が零になる。
従って全 $p_i\ne p_{i+b}$ であり、この色の反復対 $w_{i,i+b,i,i+b}$ は四つとも非零である。
もう一方の切断でも同様に色 $e$ の反復対が全て非零になる。
ゆえに色 $c$ の反復対に一つでも零があれば $(b,c),(c,b)$ の位置は、
どの $b\ne0,c$ についても四項が全て非零である。
\end{proof}

\begin{proof}[二組の文字の間の四つの主元]
色 $c$ の四反復対を全て零とする。二組のラベル $A=\{0,c\}$, $B=K\setminus A$ を取り、
$A$ の二文字を先に、$B$ の二文字を後に並べる。色 $c$ の二行二列の対角ブロックは
対向反復対と欠色により両方とも零で、非対角の二ブロックは全成分非零である。
各階数1なので $A$ 内に一つ、$B$ 内に一つ、相異なる二文字の両順を非零係数で含む関係がある。

色 $b\ne0,c$ の行列で $B=\{j,k\}$ の二行と $j+b,k+b$ の二列を取る。
対角は欠色で零、非対角は先の $(b,c)$ の全非零条件で非零である。
従って $BA$ で始まる二座標への関係平面の射影は可逆で、二つの $BA$ 語を $AB$ のみで書き換えられる。
他方の切断は $(c,b)$ を使い、反対側の列も同様である。
全四位置で六つの先頭語は降順対の全体となる。停止する二次簡約による三次元数の上界は20、
一文字縮約四本が独立で三次交叉に入ることによる下界も $64-48+4=20$ である。
縮約の単射性は異なるラベルで異なる総次数を持ち、各ラベルが実際に現れることによる。
従って三文字の重なりは全て解消し、Diamond補題で全次数の基底は非減少語となる。
二次PBWからKoszul、双対の次元列は $1,4,6,4,1$ である。一文字縮約の両側単射性と切断階数6から
最上テンソルの対合が次数1対3、2対2でも完全なので左右AS性が従う。

色 $c$ のうち一辺の両向きだけが零なら、その補辺の両ラベルを四位置とも $z$ 倍して $z^2$ で割る。
四文字が異なる項と他色の反復対は不変、補辺の両反復対だけがさらに $z^2$ 倍になる。
極限でも上の主元は残り、色 $c$ の非対角ブロックも非零で、元の階数2以下という上界は保たれる。
従って極限が上記の正則PBW窓になり、命題\ref{prop:unrestricted-good-reduction}で結論を得る。
全反復対が非零の場合は星の正規化が直接可能である。
\path{computations/verify_unmarked_klein_pbw.py} は四直線の全ての一致分割、二切断の反転と主元の色を検算する。
\end{proof}
